\section{Concentration and the successful event}\label{sec:concentration}

We first bound the magnitude of the ridge weights produced by LSVI-UCB, then
use a self-normalized concentration argument to define a high-probability event
on which least-squares value iteration tracks the Bellman backup.

\begin{lemma}[Bounded weights]\label{lem:weight-bound}
Under \Cref{ass:linear-mdp} with $\lambda=1$, for every $(k,h)\in[K]\times[H]$
the LSVI-UCB weight $w_h^k$ of Eq.~\eqref{eq:weight-def} satisfies
$\norm{w_h^k}\le 2H\sqrt{dk}$.
\end{lemma}

\begin{proof}
We bound $|v^\top w_h^k|$ for an arbitrary unit vector $v$ and take the
supremum. Writing $\featmap_h^\tau:=\featmap(x_h^\tau,a_h^\tau)$ and
$y_h^\tau:=r_h(x_h^\tau,a_h^\tau)+V_{h+1}^k(x_{h+1}^\tau)\in[0,H]$,
\begin{align*}
|v^\top w_h^k|
&\;=\; \Big| \sum_{\tau=1}^{k-1} v^\top (\Gram_h^k)^{-1}\featmap_h^\tau\, y_h^\tau \Big| \\
&\;\le\; H \sum_{\tau=1}^{k-1} \big| v^\top (\Gram_h^k)^{-1}\featmap_h^\tau \big| \\
&\;\le\; H \Big( \sum_{\tau=1}^{k-1} v^\top (\Gram_h^k)^{-1} v \Big)^{1/2} \Big( \sum_{\tau=1}^{k-1} (\featmap_h^\tau)^\top (\Gram_h^k)^{-1}\featmap_h^\tau \Big)^{1/2} \\
&\;\le\; H \cdot \norm{v} \cdot \big( dk \big)^{1/2}
\;=\; H\sqrt{dk},
\end{align*}
where the first step is Eq.~\eqref{eq:weight-def}, the second uses
$|y_h^\tau|\le H$ and the triangle inequality, the third is the Cauchy--Schwarz
inequality in the inner product $\langle u,u'\rangle_{(\Gram_h^k)^{-1}}=u^\top
(\Gram_h^k)^{-1}u'$, and the fourth uses $v^\top(\Gram_h^k)^{-1}v\le
\norm{v}^2/\lambda=\norm{v}^2$ (since $\lambda_{\min}(\Gram_h^k)\ge\lambda=1$)
together with $\sum_{\tau=1}^{k-1}(\featmap_h^\tau)^\top(\Gram_h^k)^{-1}
\featmap_h^\tau\le d$ (\Cref{lem:trace-bound}) and $\norm{v}=1$ (so the first
factor is $(\sum_{\tau=1}^{k-1}\norm{v}^2)^{1/2}=\sqrt{k-1}\le\sqrt k$). Taking the
supremum over $\norm{v}=1$ and using $\norm{w_h^k}=\sup_{\norm{v}=1}|v^\top
w_h^k|$ gives $\norm{w_h^k}\le H\sqrt{dk}\le 2H\sqrt{dk}$. This completes the
proof.
\end{proof}

The core statistical guarantee is the following self-normalized bound on the
fluctuation of the empirical Bellman backup. We define the successful event
$\Ev$ here; all later derivations are carried out on $\Ev$.

\begin{lemma}[Self-normalized concentration; successful event]\label{lem:concentration}
Suppose \Cref{ass:linear-mdp} holds and $\lambda=1$. There exist universal
constants $C_\beta>0$ and $C>0$ such that, setting
$\beta=C_\beta\,dH\sqrt{\iota}$ with $\iota=\log(2dT/\delta)$, the event
\begin{align}\label{eq:event-def}
\Ev \;:=\; \bigg\{
\forall (k,h)\in[K]\times[H]:\;
\Big\| \sum_{\tau=1}^{k-1} \featmap_h^\tau \big[\, V_{h+1}^k(x_{h+1}^\tau) - (\Pop_h V_{h+1}^k)(x_h^\tau,a_h^\tau) \,\big] \Big\|_{(\Gram_h^k)^{-1}}
\;\le\; C\, dH\sqrt{\iota}
\bigg\}.
\end{align}
satisfies $\Prob[\Ev]\ge 1-\delta/2$.
\end{lemma}

\begin{proof}
The argument is the self-normalized tail bound for vector-valued martingales of
\cite{abbasi2011improved}, combined with a covering-number union bound over the
class of value functions $V_{h+1}^k$ can equal; we record its structure and the
constant accounting.

Fix $(k,h)$ and abbreviate $\featmap_h^\tau:=\featmap(x_h^\tau,a_h^\tau)$. For a
\emph{fixed} function $V:\St\to[0,H]$, the increments
$\varepsilon_h^\tau(V):= V(x_{h+1}^\tau) - (\Pop_h V)(x_h^\tau,a_h^\tau)$ form a
martingale-difference sequence with respect to $\{\mathcal F_{\tau,h}\}$ with
$\E[\varepsilon_h^\tau(V)\mid\mathcal F_{\tau,h-1}]=0$ and
$|\varepsilon_h^\tau(V)|\le H$, hence conditionally $H$-sub-Gaussian; the
regressors $\featmap_h^\tau$ are $\mathcal F_{\tau,h-1}$-measurable, hence
$\{\mathcal F_{\tau,h-1}\}$-predictable, so the hypotheses of
\Cref{lem:self-normalized} are met. By the
self-normalized bound of \cite{abbasi2011improved} (recorded in
\Cref{lem:self-normalized}), for any fixed $V$ and any $\delta'\in(0,1)$, with
probability at least $1-\delta'$,
\begin{align}
\Big\| \sum_{\tau=1}^{k-1} \featmap_h^\tau\, \varepsilon_h^\tau(V) \Big\|_{(\Gram_h^k)^{-1}}^2
&\;\le\; 2H^2 \log\!\bigg( \frac{\det(\Gram_h^k)^{1/2}\det(\lambda I)^{-1/2}}{\delta'} \bigg) \notag \\
&\;\le\; 2H^2 \Big( \tfrac{d}{2}\log\tfrac{\lambda+k}{\lambda} + \log\tfrac{1}{\delta'} \Big), \label{eq:sn-fixed}
\end{align}
where the first step is \Cref{lem:self-normalized} with $\sigma=H$ and the
second bounds $\det(\Gram_h^k)\le(\lambda+k)^d$ (the eigenvalues of $\Gram_h^k$
sum to at most $\lambda d+k$, so their product is at most $(\lambda+k/d)^d\le
(\lambda+k)^d$ by AM--GM) and $\det(\lambda I)=\lambda^d$.

The function $V_{h+1}^k$ in Eq.~\eqref{eq:event-def} is \emph{not} fixed: it
depends on the data through $(w_{h+1}^k,\Gram_{h+1}^k)$. By
Eq.~\eqref{eq:Q-def} and Eq.~\eqref{eq:V-def} together with
\Cref{lem:weight-bound}, every such
$V_{h+1}^k$ lies in the parametric class
\begin{align*}
\Vsp \;:=\; \Big\{\, x\mapsto \min\big\{\max_a \big[\, \inner{w}{\featmap(x,a)} + \beta\,\Mnorm{\featmap(x,a)}{A^{-1}} \,\big],\,H\big\} : \norm{w}\le 2H\sqrt{dK},\; A\succeq I \,\Big\}.
\end{align*}
Let $\Vsp_\epsilon$ be a minimal $\epsilon$-cover of $\Vsp$ in the supremum
norm. A volumetric estimate over the parameters $(w,A)$ (\cite{jin2020provably},
Lemma D.6) gives
\begin{align}\label{eq:cover}
\log|\Vsp_\epsilon| \;\le\; C_1\, d^2 \log\!\big( 1 + \tfrac{HdK}{\epsilon\lambda} \big) .
\end{align}
Apply Eq.~\eqref{eq:sn-fixed} with $\delta'=\delta/(2HK|\Vsp_\epsilon|)$ and
union-bound over $\Vsp_\epsilon$ and over all $(k,h)$, so that the fixed-$V$
bound holds simultaneously for every $\bar V\in\Vsp_\epsilon$. The realized
$V_{h+1}^k\in\Vsp$ need not lie in $\Vsp_\epsilon$, so we transfer to its nearest
cover element $\bar V$ with $\norm{V_{h+1}^k-\bar V}_\infty\le\epsilon$: since
$\varepsilon_h^\tau$ is $1$-Lipschitz in $V$ under $\norm{\cdot}_\infty$ and
$\Mnorm{\sum_\tau\featmap_h^\tau\cdot}{(\Gram_h^k)^{-1}}$ is a seminorm, the
discretization correction is additive,
\begin{align*}
\Big\| \sum_{\tau=1}^{k-1} \featmap_h^\tau\, \varepsilon_h^\tau(V_{h+1}^k) \Big\|_{(\Gram_h^k)^{-1}}
\;\le\; \Big\| \sum_{\tau=1}^{k-1} \featmap_h^\tau\, \varepsilon_h^\tau(\bar V) \Big\|_{(\Gram_h^k)^{-1}}
\;+\; 2\epsilon\sqrt{k}\,,
\end{align*}
the second term bounding $\Mnorm{\sum_\tau\featmap_h^\tau\,[\varepsilon_h^\tau
(V_{h+1}^k)-\varepsilon_h^\tau(\bar V)]}{(\Gram_h^k)^{-1}}\le 2\epsilon
(\sum_\tau(\featmap_h^\tau)^\top(\Gram_h^k)^{-1}\featmap_h^\tau)^{1/2}\le
2\epsilon\sqrt{d}\le 2\epsilon\sqrt k$ via Cauchy--Schwarz and
\Cref{lem:trace-bound}. With $\epsilon=dH/K$ this correction is at most
$2dH\sqrt K/K\le 2dH$, lower order than the radius below. Hence, with
probability at least $1-\delta/2$, simultaneously for all $(k,h)$,
\begin{align*}
\Big\| \sum_{\tau=1}^{k-1} \featmap_h^\tau\, \varepsilon_h^\tau(V_{h+1}^k) \Big\|_{(\Gram_h^k)^{-1}}^2
&\;\le\; 2\cdot 2H^2\Big( \tfrac{d}{2}\log(\lambda+k) + \log\tfrac{1}{\delta'} \Big) + 2(2dH)^2 \\
&\;\le\; 2H^2\Big( \tfrac{d}{2}\,\iota + C_2\, d^2\,\iota \Big) \\
&\;\le\; C^2 d^2 H^2 \iota,
\end{align*}
where the first inequality applies Eq.~\eqref{eq:sn-fixed} to $\bar V$ together
with the additive correction above (using $(a+b)^2\le2a^2+2b^2$), the second
substitutes
$\log(1/\delta')\le C_2 d^2\iota$ from Eq.~\eqref{eq:cover} (with $\iota=
\log(2dT/\delta)$ absorbing the polynomial arguments $\log(1+HdK/(\epsilon
\lambda))\lesssim\iota$), and the last collects the dominant $d^2$ term into a
single universal $C^2$. Taking square roots yields Eq.~\eqref{eq:event-def} with
radius $C\,dH\sqrt{\iota}$, where $C$ denotes this fixed event-radius constant
(the same $C$ that reappears in \Cref{lem:recursion}); choosing $C_\beta\ge C+4$
in $\beta=C_\beta dH\sqrt{\iota}$ (the stronger of the two floors $C_\beta\ge C$
here and $C_\beta\ge C+4$ in \Cref{lem:recursion}) makes $\beta$ dominate this
fluctuation. The constant $C_\beta$ is thus a universal constant---independent of
$d,H,K,T,\delta$---whose existence and $dH\sqrt{\iota}$ scaling are exactly the
exploration-bonus normalization of \cite{jin2020provably} (their Theorem~3.1,
$\beta=c\,dH\sqrt{\iota}$): the constant is fixed only implicitly, by the
self-consistency requirement that $\beta$ exceed the event radius $C\,dH\sqrt{\iota}$
(equivalently their union-bound balance $c'\sqrt{\log 2+\log(c_\beta+1)}\le
c_\beta\sqrt{\log 2}$), and admits no closed-form numeric value; only the rate
factors $dH\sqrt{\iota}$ are pinned. Hence $\Prob[\Ev]\ge 1-\delta/2$. This
completes the proof.
\end{proof}
